\booksection{Annexe : analyse de l'extension de mémoire Android}

\chapterintro{Lors de l'achat d'un smartphone, vous voyez souvent la publicité « 12 Go + 8 Go RAM » — comme si 8 Go de RAM apparaissaient par magie. Mais cela rend-il vraiment le smartphone plus rapide ? Et ne risque-t-il pas d'abîmer le smartphone ? Dans cette section, nous allons expliquer à fond la technologie « d'extension de mémoire ».}

\booksubsection{I. Définition essentielle}

L'« extension de mémoire » Android (aussi appelée extension de RAM, fusion de mémoire, RAM Plus, etc.)\textbf{n'ajoute pas physiquement de puce de RAM}, mais est une technologie logicielle au niveau système : elle convertit dynamiquement une partie du stockage interne (ROM/flash UFS) pour simuler une RAM virtuelle.

Pour ainsi dire : la RAM physique est le « plan de travail rapide », la mémoire étendue est le « débarras temporaire ». Quand le plan de travail est insuffisant, le système déplace les données inutilisées vers le débarras, et les ramène quand nécessaire. Le débarras est grand, mais l'aller-retour est forcément plus lent que de laisser sur le plan de travail.

\booksubsection{II. Principe technique}

Les principaux fabricants l'implémentent principalement de deux façons :

\subsubsection*{Compression zRAM}

Divise un espace de données compressées dans la RAM, compresse le code d'applications inactives, taux de compression d'environ 2:1. Les données restent dans la RAM, rapides, sans user la flash, mais le CPU participe à la compression/décompression.

\subsubsection*{Partition swap}

Divise directement un espace fixe dans la flash UFS comme mémoire virtuelle, écrit les données des processus en arrière-plan dans le stockage. Dépend de la bande passante de stockage, vitesse de lecture/écriture bien inférieure à la mémoire physique.

La publicité « 12 Go + 8 Go RAM » est généralement une combinaison de 12 Go de RAM physique + 8 Go de mémoire virtuelle flash. À noter, zRAM n'occupe pas par défaut le stockage interne ; ce n'est que lorsque zRAM est complètement rempli de données que le système crée un fichier d'échange sur le stockage, et cette situation est extrêmement rare.

\booksubsection{III. L'écart de vitesse est le goulot clé}

La limite fondamentale de cette technologie réside dans l'écart de vitesse. La flash UFS 4.0 actuelle a une vitesse de lecture séquentielle d'environ 4200 Mo/s, soit seulement 1/10 à 1/15 de la mémoire physique LPDDR5X. Ainsi, pendant la phase de compression zRAM à faible charge, la vitesse est proche de la RAM réelle, mais dès qu'on entre dans des scénarios de forte charge, la vitesse chute à quelques dizaines de fois moins.

\booksubsection{IV. Scénarios applicables et limites des bénéfices}

L'extension de mémoire n'est pas universelle, l'effet varie selon le modèle et le scénario d'utilisation :

\begin{itemize}
  \item \textbf{Cas utiles}: modèles avec 6 Go de RAM physique + UFS 3.1 ou supérieur et espace libre $\geq$ 30 Go ; après activation d'une extension de 3 Go, le nombre d'apps maintenues en arrière-plan augmente de 3-4, le taux de rechargement baisse de 68 \%, adapté aux livreurs et chauffeurs VTC, utilisateurs intensifs de multitâche.
  \item \textbf{Scénarios quotidiens}: amélioration notable pour le multi-comptes WeChat, le maintien des onglets de navigateur, la réactivité des messages de messagerie instantanée, le temps de démarrage à froid des apps en arrière-plan considérablement réduit.
  \item \textbf{Scénarios limités}: amélioration limitée pour le chargement de gros jeux, le montage vidéo 4K et autres scénarios à forte bande passante.
  \item \textbf{Cas nuisibles}: modèles avec 8 Go de RAM + UFS 2.2 et espace libre insuffisant ; après activation, les fluctuations de fréquence d'images du jeu s'aggravent, la consultation de courtes vidéos consomme plus de batterie, le basculement saccade.
  \item \textbf{Cas inutiles}: les modèles avec plus de 16 Go de RAM n'en tirent presque aucun bénéfice, maintiennent au mieux une app supplémentaire en arrière-plan, mais augmentent la charge CPU de 12 \% et la chaleur de 2-3 °C. Pour les modèles avec plus de 12 Go de RAM, il est recommandé de désactiver pour éviter des calculs de compression inutiles.
\end{itemize}

\booksubsection{V. Dénomination et activation par fabricant}

Les marques l'appellent différemment : Xiaomi « extension de mémoire », Huawei « fusion de mémoire », Samsung « RAM Plus », vivo « Extended RAM », OnePlus « Memory Extension ». Cette fonction est une capacité native du système, sans lien direct avec la version sous-jacente d'Android.

Le chemin d'activation général : [Réglages] [Plus de réglages] trouver la sous-section [Mémoire] ou [Performance] activer « extension de mémoire » et choisir le palier de capacité (généralement 2 Go-7 Go), redémarrer pour prendre effet. Après activation, le système affiche une combinaison de mémoire telle que « 8 Go + 4 Go ».

\begin{tcbwarn}{Précautions}
\begin{itemize}
  \item L'extension de mémoire occupe une partie du stockage interne, veillez à ce qu'il reste au moins 20-30 Go, sinon la fluidité du système peut être affectée.
  \item La flash utilisée comme mémoire virtuelle subit une usure de lecture/écriture ; bien que la flash UFS moderne ait une longue durée de vie, un swap intense à long terme peut accélérer le vieillissement du stockage.
  \item Si la RAM physique du smartphone est faible (4 Go ou moins), l'extension de mémoire ne fait que soulager une urgence, sans résoudre fondamentalement le problème de lag.
\end{itemize}

\end{tcbwarn}

\begin{tcbinfo}{Résumé}
L'extension de mémoire est une technologie d'optimisation au niveau système mature, qui peut améliorer les capacités de multitâche et le maintien en arrière-plan des modèles milieu et entrée de gamme sans modifier le matériel, mais elle est limitée par le fossé de vitesse entre flash et RAM physique, et n'est pas une « panacée ». L'activer ou non, et à quel niveau, dépend de la taille de la RAM physique, de la spécification de flash, de l'espace de stockage libre et des habitudes d'utilisation.

\end{tcbinfo}

\begin{tcbpractice}{À vous de jouer : vérifiez si votre smartphone a l'extension de mémoire}
\begin{enumerate}
  \item Ouvrez [Réglages], recherchez « extension de mémoire », « fusion de mémoire », « RAM Plus » ou « Extended RAM ».
  \item Une fois trouvé, vérifiez si c'est activé actuellement et quel est le réglage de capacité étendue.
  \item Selon cette section, évaluez : quelle est la taille de RAM physique de votre modèle ? Quelle est la spécification de flash ? L'espace de stockage libre est-il suffisant ?
  \item Essayez de désactiver ou de réduire la capacité étendue, observez pendant une semaine l'évolution du maintien en arrière-plan et de la chauffe, pour trouver le réglage qui vous convient le mieux.
\end{enumerate}

\end{tcbpractice}
